\section{Registry-Struktur}

Die Windows Registry besteht aus zwei Hauptkategorien:

\begin{description}
    \item[HKLM (HKEY\_LOCAL\_MACHINE)] — Systemweite Einstellungen
    \begin{itemize}
        \item \textbf{SAM} — Benutzerkonten, Passwort-Hashes
        \item \textbf{SYSTEM} — Hardwarekonfiguration, USB-Historie, Boot-Konfiguration
        \item \textbf{SOFTWARE} — Installierte Software, OS-Version
        \item Speicherort: \texttt{Windows/System32/config/}
    \end{itemize}
    \item[HKCU (HKEY\_CURRENT\_USER)] — Benutzerspezifische Einstellungen
    \begin{itemize}
        \item \textbf{NTUSER.DAT} — Benutzereinstellungen, zuletzt geöffnete Dokumente
        \item Speicherort: \texttt{Users/<Benutzername>/NTUSER.DAT}
    \end{itemize}
\end{description}

\section{Hives extrahieren}

\begin{lstlisting}[language=bash]
fls -pr /dev/loopXpN | grep 'SAM$'
icat /dev/loopXpN <INODE_SAM> > /tmp/SAM

fls -pr /dev/loopXpN | grep 'SYSTEM$'
icat /dev/loopXpN <INODE_SYSTEM> > /tmp/SYSTEM

fls -pr /dev/loopXpN | grep 'SOFTWARE$'
icat /dev/loopXpN <INODE_SOFTWARE> > /tmp/SOFTWARE

fls -pr /dev/loopXpN | grep 'NTUSER.DAT$'
icat /dev/loopXpN <INODE_NTUSER> > /tmp/NTUSER.DAT
\end{lstlisting}

Nach jeder Extraktion mit \texttt{file /tmp/HIVE} verifizieren.

\section{Regripper-Plugins}

\begin{center}
\begin{tabular}{lll}
    \toprule
    \textbf{Plugin} & \textbf{Hive} & \textbf{Information} \\
    \midrule
    winver & SOFTWARE & Windows-Version, Build-Nr., Install-Datum \\
    usbdevices & SYSTEM & USB-Gerätenamen, VID/PID \\
    usbstor & SYSTEM & USB-Timestamps (first/last), Seriennummern \\
    usbstor2 & SYSTEM & Erweiterte USB-Timestamps \\
    devclass & SYSTEM & Geräteklassen inkl. Hersteller \\
    mountdev & SYSTEM & Laufwerksbuchstaben der USB-Geräte \\
    nic2 & SYSTEM & Netzwerkadapter, DHCP-Server, IP-Adressen \\
    run & SOFTWARE / NTUSER.DAT & Autostart-Einträge \\
    recentdocs & NTUSER.DAT & Zuletzt geöffnete Dokumente \\
    shellbags & NTUSER.DAT & Geöffnete Ordner (Explorer-Verlauf) \\
    uninstall & SOFTWARE & Installierte/deinstallierte Programme \\
    \bottomrule
\end{tabular}
\end{center}

\begin{lstlisting}[language=bash]
# Plugin suchen
regripper -l | grep -B 1 USB

# Plugin ausfuehren
regripper -p winver -r SOFTWARE
regripper -p usbdevices -r SYSTEM
regripper -p usbstor -r SYSTEM
regripper -p nic2 -r SYSTEM
\end{lstlisting}

\textbf{Hinweis:} Windows 11 meldet sich in der Registry weiterhin als ``Windows 10'' — Microsoft hat die Versionsnummer nicht mehr erhöht.

\section{Windows-Login umgehen}

\begin{description}
    \item[Utility Manager Trick] \texttt{utilman.exe} durch \texttt{cmd.exe} ersetzen.
    Beim Login-Screen den Utility Manager öffnen → Shell-Zugang.
    \item[chntpw] Interaktives Passwort-Reset über die SAM-Datei:
\end{description}

\begin{lstlisting}[language=bash]
chntpw -i /tmp/SAM
\end{lstlisting}
